\documentclass[11pt,a4paper]{article}

\usepackage{estilo-libro}

% ---------- Macros propias de este documento ----------
\newcommand{\Sini}{\ensuremath{S^{*}}}
\newcommand{\Tfin}{\ensuremath{T^{*}}}
\newcommand{\x}[2]{x_{\mathrm{#1}\,\mathrm{#2}}}
\newcommand{\ra}{\ensuremath{\to}}
\tikzset{
  regreso/.style={-{Stealth[length=2.6mm]}, line width=2.2pt, dash pattern=on 5pt off 2.5pt},
  gris/.style={-{Stealth[length=2mm]}, thin, gray!55},
  calle/.style={thick},
  ofe/.style={font=\scriptsize},
}

\title{\textbf{El problema de flujo a costo mínimo}\\[4pt]\large Explicación del método y ejercicios resueltos}
\author{Investigación de Operaciones · Asignación y flujo en redes}
\date{}

\begin{document}
\maketitle
\tableofcontents
\newpage


% =====================================================================
\part{¿En qué consiste el problema de flujo a costo mínimo?}

\section{¿Qué tipo de problema resuelve?}

Muchas empresas tienen que \textbf{mover cosas} desde los lugares donde se producen hasta los lugares donde se
necesitan: mercadería, agua, electricidad, pasajeros, mensajes. Casi siempre hay varios caminos posibles, cada
uno con un precio distinto, y algunos caminos tienen un \textbf{límite} de lo que pueden llevar. La pregunta es
siempre la misma: \textbf{¿por dónde mandar cada cosa para entregar todo lo que se pide gastando lo menos posible?}
A eso se le llama el \textbf{problema de flujo a costo mínimo}.

Para entenderlo usaremos un ejemplo pequeño que acompañará toda la Parte I:

\begin{ojo}
Dos panaderías, \textbf{A} y \textbf{B}, hornean cada mañana \textbf{10 cajas de pan} cada una. Dos tiendas,
\textbf{X} y \textbf{Y}, piden \textbf{10 cajas} cada una. Llevar una caja cuesta (en soles):
de A a X, 1; de A a Y, 3; de B a X, 2; de B a Y, 8. ¿Qué panadería debe abastecer a qué tienda para que el
reparto salga lo más barato posible?
\end{ojo}

La primera idea que se le ocurre a cualquiera es ``usar primero lo más barato'': A le manda sus 10 cajas a X
(cuesta 1 cada una). Pero entonces a B solo le queda la tienda Y, que es su ruta más cara (8 por caja), y el
reparto cuesta $10 + 80 = 90$. Veremos que existe un plan mucho mejor y que el método lo encuentra solo,
\textbf{deshaciendo} la primera decisión cuando conviene.

Todo problema de flujo a costo mínimo se dibuja como una \textbf{red}: puntos (\textbf{nodos}) unidos por flechas
(\textbf{arcos}). Hay nodos de \textbf{oferta}, que producen y llevan un número positivo (A y B: $+10$); nodos de
\textbf{demanda}, que necesitan y llevan un número negativo (X y Y: $-10$); y, a veces, nodos de \textbf{transbordo},
puntos de paso como un almacén, que no producen ni piden (0). Cada arco tiene un \textbf{costo por unidad} y, a veces,
una \textbf{capacidad}. Lo que hay que decidir es el \textbf{flujo}: cuánto se manda por cada arco.

% =====================================================================
\section{Las reglas del problema}

Un plan de envíos solo es válido si cumple dos reglas:

\begin{description}[leftmargin=0pt, itemsep=5pt]
  \item[Conservación (nada aparece ni desaparece).] En cada nodo,
        \[ \text{lo que sale} - \text{lo que entra} = \text{su oferta}. \]
        En una panadería sale lo que hornea ($+10$); en una tienda entra lo que pide (sale $-$ entra $= -10$);
        en un nodo de transbordo sale exactamente lo que entra ($= 0$).
  \item[Capacidad.] Por cada arco no se puede mandar más de su capacidad, ni cantidades negativas.
\end{description}

Entre todos los planes que cumplen las reglas se busca el de \textbf{menor costo total}, que es la suma de
(costo por unidad) $\times$ (unidades enviadas) en todos los arcos. Escrito como programación lineal, con
$x_{ij}$ = unidades que van del nodo $i$ al nodo $j$, $c_{ij}$ su costo, $u_{ij}$ su capacidad y $b_i$ la oferta del nodo $i$:
\[
\begin{aligned}
\text{Min } Z &= \sum_{\text{arcos}} c_{ij}\, x_{ij}\\
\text{s.a.}\quad \sum_{j} x_{ij} - \sum_{k} x_{ki} &= b_i \qquad \text{para cada nodo } i \quad (\text{sale} - \text{entra})\\
0 \le x_{ij} &\le u_{ij} \qquad \text{para cada arco}
\end{aligned}
\]
En el ejemplo de las panaderías:
\[
\begin{aligned}
\text{Min } Z &= 1\,\x{A}{X} + 3\,\x{A}{Y} + 2\,\x{B}{X} + 8\,\x{B}{Y}\\
\x{A}{X} + \x{A}{Y} &= 10 \nota{lo que hornea A}\\
\x{B}{X} + \x{B}{Y} &= 10 \nota{lo que hornea B}\\
-\x{A}{X} - \x{B}{X} &= -10 \nota{lo que pide X}\\
-\x{A}{Y} - \x{B}{Y} &= -10 \nota{lo que pide Y}
\end{aligned}
\]

Este modelo se podría resolver con el Simplex (por ejemplo, el Simplex Revisado del primer previo), pero la forma
de red permite un método mucho más directo, que se explica a continuación. Hay además un regalo:
\textbf{si las ofertas, las demandas y las capacidades son números enteros, la mejor solución también sale con
números enteros}. Nunca hará falta mandar ``media caja''.

% =====================================================================
\section{El problema ``padre'': sus casos particulares}

El flujo a costo mínimo es el problema más general de los que se ven en este curso: muchos problemas conocidos
son el mismo problema con datos especiales. Por eso, un solo método sirve para todos ellos.

\begin{center}\small
\begin{tabular}{p{2.9cm}p{7.1cm}p{4.2cm}}
\toprule
\textbf{Problema} & \textbf{Cómo se ve como flujo a costo mínimo} & \textbf{En este documento} \\
\midrule
Transporte & Solo nodos de oferta y de demanda; un arco directo de cada origen a cada destino, sin capacidad. & Ejercicio 1 (Powerco) \\
Asignación & Un transporte donde cada persona ofrece 1 y cada tarea pide 1. & tema de asignación \\
Transbordo & Hay nodos de paso (oferta 0) por los que la mercadería puede pasar antes de llegar. & Ejercicio 2 \\
Ruta más corta & El origen ofrece 1 unidad, el destino pide 1, y el costo de cada arco es su distancia. La
                 ruta por donde viaja esa unidad es la más corta. & Ejercicio 3 (Seervada) \\
Flujo máximo & Los arcos cuestan 0 y se agrega un atajo de la fuente al destino carísimo y sin límite. Se pide mandar
               una cantidad enorme: para gastar poco, el método mete por la red todo lo que cabe (el flujo máximo). & tema de flujo máximo \\
\bottomrule
\end{tabular}
\end{center}

% =====================================================================
\section{La idea del método: rutas más baratas sucesivas}

El método que usamos se llama \textbf{rutas más baratas sucesivas} (en los libros, \emph{caminos más cortos
sucesivos}). La idea es la de un despachador sensato:

\begin{enumerate}[itemsep=3pt]
  \item Busca, entre todos los caminos que todavía tienen espacio, \textbf{la ruta más barata} desde algún punto que
        aún tiene producto hasta algún punto que aún lo necesita.
  \item Manda por ella \textbf{todo lo que se pueda}, hasta que algo se agote: la oferta que queda, lo que falta
        por entregar o el espacio de algún tramo.
  \item Repite hasta que \textbf{todo lo pedido esté entregado}.
\end{enumerate}

Lo que hace que el método no se equivoque es un detalle: \textbf{también puede deshacer envíos anteriores}. Si
antes se mandó algo por un arco, ahora se permite recorrer ese arco al revés, y hacerlo \textbf{resta} su costo,
porque significa ``ya no mando eso por ahí''. En el ejemplo:

\begin{center}
\begin{tikzpicture}[x=1.15cm, y=1.15cm]
  \begin{scope}
    \node[fuente, label=left:{\small $+10$}] (A) at (0,0) {A};
    \node[fuente, label=left:{\small $+10$}] (B) at (0,-2) {B};
    \node[sumidero, label=right:{\small $-10$}] (X) at (3,0) {X};
    \node[sumidero, label=right:{\small $-10$}] (Y) at (3,-2) {Y};
    \draw[arco] (A) -- node[etq, above] {1} (X);
    \draw[arco] (A) -- node[etq, pos=0.25] {3} (Y);
    \draw[arco] (B) -- node[etq, pos=0.25] {2} (X);
    \draw[arco] (B) -- node[etq, below] {8} (Y);
    \node[font=\small\bfseries] at (1.5,-3) {La red (costo por caja)};
  \end{scope}
  \begin{scope}[xshift=5.6cm]
    \node[fuente] (A) at (0,0) {A};
    \node[fuente] (B) at (0,-2) {B};
    \node[sumidero] (X) at (3,0) {X};
    \node[sumidero] (Y) at (3,-2) {Y};
    \draw[camino] (A) -- node[etq, above] {$+1$} (X);
    \draw[gris] (A) -- (Y); \draw[gris] (B) -- (X); \draw[gris] (B) -- (Y);
    \node[font=\small\bfseries, align=center] at (1.5,-3.2) {Ruta 1: A $\to$ X\\[-1pt]\footnotesize 10 cajas a 1 = 10};
  \end{scope}
  \begin{scope}[xshift=11.2cm]
    \node[fuente] (A) at (0,0) {A};
    \node[fuente] (B) at (0,-2) {B};
    \node[sumidero] (X) at (3,0) {X};
    \node[sumidero] (Y) at (3,-2) {Y};
    \draw[gris] (B) -- (Y);
    \draw[camino] (B) -- node[etq, pos=0.22] {$+2$} (X);
    \draw[regreso] (X) -- node[etq, above] {$-1$} (A);
    \draw[camino] (A) -- node[etq, pos=0.22] {$+3$} (Y);
    \node[font=\small\bfseries, align=center] at (1.5,-3.2) {Ruta 2: B $\to$ X $\to$ A $\to$ Y\\[-1pt]\footnotesize 10 cajas a $2-1+3 = 4$ = 40};
  \end{scope}
\end{tikzpicture}
\end{center}

La ruta 2 se lee así: B le manda 10 cajas a X (cuesta 2 cada una); como X ya tenía las cajas de A, A deja de
mandárselas (ahorra 1 cada una, por eso el tramo punteado \textbf{al revés} vale $-1$); y A usa esas 10 cajas para
abastecer a Y (cuesta 3). Al final A abastece a Y y B a X: el reparto cuesta $10 + 40 = 50$, en lugar de 90.
Nadie mandó cajas ``hacia atrás'': simplemente \textbf{se corrigió el plan}.

% =====================================================================
\section{Antes de empezar: preparar el problema}

\begin{description}[leftmargin=0pt, itemsep=5pt]
  \item[Revisar si está balanceado.] Se suma lo que se ofrece y lo que se pide. Si son iguales, se sigue.
        Si se ofrece más de lo que se pide, se agrega un \textbf{destino ficticio} que ``pide'' el sobrante, con costo 0
        desde cada origen: lo que se ``manda'' ahí es lo que en realidad se queda en casa. Si se pide más de lo que se
        ofrece, se agrega un \textbf{origen ficticio} que ``ofrece'' lo que falta: lo que sale de él es demanda que se
        queda sin atender.
  \item[Juntar las ofertas y las demandas.] Para buscar una sola ruta a la vez, se agrega un \textbf{punto inicial
        \Sini} con un arco de costo 0 hacia cada nodo de oferta, cuyo espacio es lo que ese nodo ofrece; y un
        \textbf{punto final \Tfin} con un arco de costo 0 desde cada nodo de demanda, cuyo espacio es lo que ese nodo
        pide. Así, ``ir de \Sini{} a \Tfin'' significa ``llevar algo de quien tiene a quien necesita''.
  \item[Arcos en los dos sentidos.] Si una conexión se puede usar en ambos sentidos (una calle de doble vía), se
        dibujan dos arcos, uno de ida y uno de vuelta, con el mismo costo.
\end{description}

En el ejemplo, se ofrecen $10 + 10 = 20$ cajas y se piden $10 + 10 = 20$: está balanceado. La red preparada queda:

\begin{center}
\begin{tikzpicture}[x=1.4cm, y=0.8cm]
  \node[nodo] (S) at (0,-1) {\Sini};
  \node[fuente] (A) at (2,0) {A};
  \node[fuente] (B) at (2,-2) {B};
  \node[sumidero] (X) at (5,0) {X};
  \node[sumidero] (Y) at (5,-2) {Y};
  \node[nodo] (T) at (7,-1) {\Tfin};
  \draw[arco, dashed] (S) -- node[etq, above left] {0; hasta 10} (A);
  \draw[arco, dashed] (S) -- node[etq, below left] {0; hasta 10} (B);
  \draw[arco] (A) -- node[etq, above] {1} (X);
  \draw[arco] (A) -- node[etq, pos=0.25] {3} (Y);
  \draw[arco] (B) -- node[etq, pos=0.25] {2} (X);
  \draw[arco] (B) -- node[etq, below] {8} (Y);
  \draw[arco, dashed] (X) -- node[etq, above right] {0; hasta 10} (T);
  \draw[arco, dashed] (Y) -- node[etq, below right] {0; hasta 10} (T);
\end{tikzpicture}
\end{center}

% =====================================================================
\section{Las palabras que se usan}

\begin{center}\small
\begin{tabular}{p{3.6cm}p{10.6cm}}
\toprule
\textbf{Término} & \textbf{Qué significa} \\
\midrule
Oferta $b_i$ & Lo que produce un nodo (positivo). Si es negativo, es lo que pide (demanda). Si es 0, es de transbordo. \\
Costo por unidad $c_{ij}$ & Lo que cuesta mandar una unidad por el arco de $i$ a $j$. \\
Capacidad $u_{ij}$ & Lo máximo que puede llevar un arco. Si no se indica, no tiene límite ($\infty$). \\
Balanceado & Cuando lo ofrecido es igual a lo pedido. Si no, se agrega un nodo ficticio. \\
\Sini{} y \Tfin & El punto inicial que reúne todas las ofertas y el punto final que reúne todas las demandas. \\
Red de lo que queda & La red con el \textbf{espacio libre} de cada arco en el momento actual. Los libros la llaman
                      \emph{red residual}. \\
Tramo al revés & Si un arco ya lleva flujo, se puede recorrer en sentido contrario para \textbf{deshacer} parte de ese
                 envío. Su costo es el del arco \textbf{con signo menos}, y su espacio es lo que el arco lleva. \\
Ruta más barata & El camino de \Sini{} a \Tfin{} en la red de lo que queda cuya suma de costos es la menor. \\
Tramo más limitado & El tramo de la ruta con menos espacio: decide cuánto se puede enviar (el ``cuello de botella''). \\
\bottomrule
\end{tabular}
\end{center}

% =====================================================================
\section{Los pasos del método}\label{sec:algoritmo}

Cada iteración responde cuatro preguntas, siempre en el mismo orden. Para que se entiendan, se muestra lo que
ocurre en el ejemplo de las panaderías.

\paso{Paso 1. ¿Cuánto espacio queda en cada arco?}
Se dibuja la red de lo que queda. Por cada arco de $i$ a $j$ que lleva $x_{ij}$ unidades:
\begin{itemize}[itemsep=2pt]
  \item hacia adelante quedan $u_{ij} - x_{ij}$ unidades de espacio, con costo $+c_{ij}$;
  \item si $x_{ij} > 0$, aparece un tramo al revés, de $j$ a $i$, con espacio $x_{ij}$ y costo $-c_{ij}$.
\end{itemize}
Los arcos que salen de \Sini{} tienen como espacio lo que a cada origen \textbf{le queda} por mandar, y los que llegan a
\Tfin{}, lo que a cada destino \textbf{le falta} recibir.

\textit{En el ejemplo:} al inicio no se ha mandado nada, así que la red de lo que queda es la red original:
a A y a B les quedan 10 cajas, a X y a Y les faltan 10.

\paso{Paso 2. ¿Cuál es la ruta más barata?}
Se busca la ruta de \Sini{} a \Tfin{} con menor suma de costos, usando solo tramos con espacio. Como los tramos al
revés tienen costo negativo, se usa un procedimiento que los admite (el de \textbf{Bellman-Ford}): se anota 0 en
\Sini{} y ``infinito'' en los demás puntos; se revisan todos los tramos y, si llegar a $i$ cuesta $d_i$ y el tramo de
$i$ a $j$ cuesta $c$, entonces a $j$ se puede llegar con $d_i + c$; si eso es menor que lo anotado en $j$, se
corrige. Se repite la revisión hasta que ya nada cambie. Lo anotado en \Tfin{} es el costo de la ruta más barata.

\textit{En el ejemplo:} las cuatro rutas posibles cuestan 1, 3, 2 y 8. La más barata es \Sini{} \ra{} A \ra{} X \ra{} \Tfin,
con costo 1 por caja.

\paso{Paso 3. ¿Cuánto se envía?}
Lo que permita el tramo más limitado de la ruta: el \textbf{menor espacio} entre todos sus tramos (incluidos lo que le
queda al origen y lo que le falta al destino).

\textit{En el ejemplo:} $\min(10,\ \infty,\ 10) = 10$: a A le quedan 10, el arco A \ra{} X no tiene límite y a X le
faltan 10. Se mandan 10 cajas.

\paso{Paso 4. Enviar y anotar}
Se suma el envío al flujo de cada arco de la ruta (y, si un tramo iba al revés, se \textbf{resta} del arco original).
El costo de la iteración es (costo por unidad) $\times$ (envío) y se suma al acumulado. Si todavía falta entregar algo,
se vuelve al paso 1.

\textit{En el ejemplo:} A \ra{} X lleva 10 cajas; costo $1 \times 10 = 10$, acumulado 10. Faltan las 10 cajas de Y,
así que se da otra vuelta.

\paso{Segunda vuelta del ejemplo: aquí aparece el tramo al revés}
Paso 1: a A ya no le queda nada y X ya está completa. Pero como A \ra{} X lleva 10, en la red de lo que queda
existe el tramo al revés X \ra{} A con espacio 10 y costo $-1$.

Paso 2: a B le quedan 10 y a Y le faltan 10. Hay dos rutas:
\[
\text{B} \to \text{Y}: \ 8 \qquad\qquad \text{B} \to \text{X} \to \text{A} \to \text{Y}: \ 2 - 1 + 3 = \mathbf{4}
\nota{la más barata}
\]
Paso 3: $\min(10,\ \infty,\ 10,\ \infty,\ 10) = 10$. \quad
Paso 4: costo $4 \times 10 = 40$, acumulado $10 + 40 = 50$. Al recorrer X \ra{} A al revés, el arco A \ra{} X
\textbf{baja} de 10 a 0. El plan final es A \ra{} Y: 10 y B \ra{} X: 10, con costo $3(10) + 2(10) = 50$. Ya se entregó
todo: se terminó.

\noindent
\begin{minipage}[c]{0.50\textwidth}
\begin{idea}[Resumen de los pasos]
\begin{enumerate}[leftmargin=*, itemsep=3pt]
  \item[\textbf{0.}] Preparar: balancear (nodo ficticio si hace falta) y agregar \Sini{} y \Tfin. Todos los flujos en 0.
  \item \textbf{Red de lo que queda:} espacio libre de cada arco ($+c$) y tramos al revés de lo ya enviado ($-c$).
  \item \textbf{Ruta más barata} de \Sini{} a \Tfin{} (Bellman-Ford, porque hay costos negativos).
  \item \textbf{Envío} = el menor espacio de la ruta.
  \item \textbf{Enviar y anotar:} actualizar flujos y costo acumulado. Si falta entregar algo, volver al paso 1.
\end{enumerate}
\end{idea}
\end{minipage}\hfill
\begin{minipage}[c]{0.47\textwidth}\centering
\begin{tikzpicture}[scale=0.78, transform shape, node distance=4mm and 6mm, >=Stealth,
  caja/.style={draw, rounded corners=2pt, fill=gray!8, align=center, minimum width=4.4cm, inner sep=5pt, font=\small},
  decision/.style={draw, diamond, aspect=2.4, fill=gray!18, align=center, inner sep=1pt, font=\small},
  fin/.style={draw, rounded corners=8pt, fill=gray!25, align=center, minimum width=2.2cm, inner sep=5pt, font=\small}]
  \node[caja] (ini) {Preparar: balancear,\\ agregar \Sini{} y \Tfin};
  \node[decision, below=of ini] (falta) {¿Falta entregar\\ algo?};
  \node[fin, right=of falta] (optimo) {\textbf{Solución}\\\textbf{óptima}};
  \node[caja, below=of falta] (resid) {1. Red de lo que queda\\ (con tramos al revés)};
  \node[decision, below=of resid] (hay) {¿Hay ruta de\\ \Sini{} a \Tfin?};
  \node[fin, right=of hay] (infact) {\textbf{Sin}\\\textbf{solución}};
  \node[caja, below=of hay] (envia) {2--4. Ruta más barata,\\ envío = menor espacio,\\ sumar al acumulado};
  \draw[->] (ini) -- (falta);
  \draw[->] (falta) -- node[above, font=\scriptsize] {no} (optimo);
  \draw[->] (falta) -- node[right, font=\scriptsize] {sí} (resid);
  \draw[->] (resid) -- (hay);
  \draw[->] (hay) -- node[above, font=\scriptsize] {no} (infact);
  \draw[->] (hay) -- node[right, font=\scriptsize] {sí} (envia);
  \draw[->] (envia.west) -- ++(-0.9,0) |- node[pos=0.25, left, font=\scriptsize, align=center] {siguiente\\iteración} (falta.west);
\end{tikzpicture}
\end{minipage}

% =====================================================================
\section{¿Cómo se sabe que terminó?}

\begin{description}[leftmargin=0pt, itemsep=5pt]
  \item[Solución óptima.] Ya se entregó todo lo pedido (los arcos que llegan a \Tfin{} están llenos). Como en cada
        vuelta se usó la ruta más barata disponible, el plan final es el más barato posible. Una forma de
        comprobarlo: en la red de lo que queda \textbf{no hay ningún recorrido circular con costo negativo}; si lo
        hubiera, se podría dar esa vuelta y ahorrar.
  \item[Problema sin solución.] Todavía falta entregar algo, pero no existe ninguna ruta con espacio de \Sini{} a \Tfin:
        las capacidades no alcanzan para cubrir lo pedido.
  \item[Problema sin límite.] Si hubiera un recorrido circular de costo negativo y sin límite de capacidad, se podría
        ahorrar sin fin. Con costos positivos (distancias, fletes) esto no ocurre; suele indicar un error en los datos.
  \item[Empates.] Si dos rutas cuestan lo mismo, se puede tomar cualquiera: el costo final es el mismo, aunque los
        envíos pueden quedar distintos (hay \textbf{varias soluciones óptimas}). Pasa en los ejercicios 1 y 3.
\end{description}

% =====================================================================
\section{Ventajas y desventajas}

\begin{center}\small
\begin{tabular}{p{7.2cm}p{7.2cm}}
\toprule
\textbf{Ventajas} & \textbf{Desventajas} \\
\midrule
\begin{itemize}[leftmargin=*, itemsep=2pt, topsep=2pt]
  \item Es intuitivo: siempre se usa la ruta más barata que queda.
  \item Corrige solo sus decisiones anteriores con los tramos al revés.
  \item Con datos enteros, los envíos salen enteros.
  \item Sirve para todos los casos particulares: transporte, transbordo, ruta más corta, flujo máximo.
  \item Cada iteración se puede revisar a mano en el dibujo.
\end{itemize}
&
\begin{itemize}[leftmargin=*, itemsep=2pt, topsep=2pt]
  \item Hay que redibujar la red de lo que queda en cada vuelta y no olvidar los tramos al revés.
  \item Buscar la ruta más barata con costos negativos (Bellman-Ford) es laborioso a mano en redes grandes.
  \item El número de vueltas puede crecer cuando las cantidades son grandes y hay muchos cuellos de botella.
\end{itemize}
\\
\bottomrule
\end{tabular}
\end{center}

\textbf{Otro método: el símplex de redes.} Hillier y Lieberman resuelven este problema con el \textbf{método símplex de
redes}, que es el Simplex de siempre adaptado a la forma de la red (cada solución básica es un ``árbol'' de arcos). Es
el método que usan los programas comerciales para redes enormes. Llega al \textbf{mismo costo mínimo}; lo que cambia
es el camino para llegar. Aquí usamos rutas más baratas porque cada paso se ve en el dibujo.

% =====================================================================
\part{Ejercicios resueltos}
\setcounter{section}{0}

Los tres ejercicios siguientes aplican, uno por uno, los pasos de la sección~\ref{sec:algoritmo}. Las rutas y los
números son exactamente los que muestra la calculadora (modo \textsc{5P Costo mín.}, memorias 13, 14 y 15).

% =====================================================================
\section{Powerco: el transporte como flujo a costo mínimo}
\textit{Fuente: Winston, Investigación de operaciones: aplicaciones y algoritmos, capítulo del modelo de transporte.}

\subsection*{Planteamiento del problema}
Powerco tiene \textbf{tres plantas eléctricas} que abastecen a \textbf{cuatro ciudades}. Cada planta puede generar
una cantidad de energía (en millones de kWh) y cada ciudad necesita una cantidad. Mandar un millón de kWh de una planta
a una ciudad tiene un costo (en dólares) que depende de la distancia. Como red, cada planta es un nodo de oferta, cada
ciudad un nodo de demanda y hay un arco de cada planta a cada ciudad (ninguno tiene límite).

\noindent
\begin{minipage}[c]{0.44\textwidth}\small\centering
\begin{tabular}{lccccc}
\toprule
 & \textbf{C1} & \textbf{C2} & \textbf{C3} & \textbf{C4} & \textbf{Oferta} \\
\midrule
Planta 1 & 8  & 6  & 10 & 9  & 35 \\
Planta 2 & 9  & 12 & 13 & 7  & 50 \\
Planta 3 & 14 & 9  & 16 & 5  & 40 \\
\midrule
Demanda  & 45 & 20 & 30 & 30 & \\
\bottomrule
\end{tabular}
\end{minipage}\hfill
\begin{minipage}[c]{0.54\textwidth}\centering
\begin{tikzpicture}[x=1.85cm, y=-1.3cm]
  \node[fuente, label=left:{\small $+35$}] (P1) at (0,0.5) {P1};
  \node[fuente, label=left:{\small $+50$}] (P2) at (0,2) {P2};
  \node[fuente, label=left:{\small $+40$}] (P3) at (0,3.5) {P3};
  \node[sumidero, label=right:{\small $-45$}] (C1) at (3.4,0) {C1};
  \node[sumidero, label=right:{\small $-20$}] (C2) at (3.4,1.33) {C2};
  \node[sumidero, label=right:{\small $-30$}] (C3) at (3.4,2.67) {C3};
  \node[sumidero, label=right:{\small $-30$}] (C4) at (3.4,4) {C4};
  \foreach \p/\c/\k in {P1/C1/8, P1/C2/6, P1/C3/10, P1/C4/9, P2/C1/9, P2/C2/12, P2/C3/13, P2/C4/7,
                        P3/C1/14, P3/C2/9, P3/C3/16, P3/C4/5}
    \draw[arco] (\p) -- node[etq, pos=0.26] {\k} (\c);
\end{tikzpicture}
\end{minipage}

La pregunta: \textbf{¿cuánta energía debe mandar cada planta a cada ciudad para cubrir todas las demandas al menor costo?}

\subsection*{Formulación del modelo}
Sea $x_{ij}$ = millones de kWh que la planta $i$ manda a la ciudad $j$ (12 variables, una por arco). Cada arco aporta
(costo por unidad) $\times$ (lo enviado). En la columna izquierda, lo que \textbf{sale} de cada planta es su oferta; en la
derecha, lo que \textbf{entra} a cada ciudad es su demanda (en la forma de red: sale $-$ entra $= -$demanda).
\[
\begin{aligned}
\text{Min } Z = {}& 8x_{11} + 6x_{12} + 10x_{13} + 9x_{14} + 9x_{21} + 12x_{22} + 13x_{23} + 7x_{24}\\
                  & + 14x_{31} + 9x_{32} + 16x_{33} + 5x_{34}\\[3pt]
\text{s.a.}\quad & x_{11} + x_{12} + x_{13} + x_{14} = 35 && x_{11} + x_{21} + x_{31} = 45\\
& x_{21} + x_{22} + x_{23} + x_{24} = 50 && x_{12} + x_{22} + x_{32} = 20\\
& x_{31} + x_{32} + x_{33} + x_{34} = 40 && x_{13} + x_{23} + x_{33} = 30\\
& x_{ij} \ge 0 && x_{14} + x_{24} + x_{34} = 30
\end{aligned}
\]
Este es el clásico \textbf{modelo de transporte}: un flujo a costo mínimo sin nodos de paso y sin capacidades.

\subsection*{Preparación}
Se ofrecen $35 + 50 + 40 = 125$ y se piden $45 + 20 + 30 + 30 = 125$: \textbf{está balanceado}, no hace falta nodo
ficticio. Se agrega \Sini{} con arcos a P1, P2 y P3 (espacio 35, 50 y 40) y \Tfin{} con arcos desde C1, \dots, C4
(espacio 45, 20, 30 y 30), todos de costo 0. Al inicio nada se ha enviado: costo acumulado 0.

% ---------------------------------------------------------------------
\subsection*{Las rutas, una por una}

En las cuatro primeras vueltas la ruta más barata es un arco directo. Dentro del mínimo van, en orden, lo que le
queda a la planta, el espacio del arco ($\infty$: no tiene límite) y lo que le falta a la ciudad.

\noindent
\begin{minipage}[t]{0.49\textwidth}
\paso{Ruta 1: \ P3 $\to$ C4}
$\begin{aligned}
\text{por unidad} &= 5 \nota{el arco más barato}\\
\text{envío} &= \min(40, \infty, 30) = 30 \nota{C4 llena}\\
\text{costo} &= 5 \times 30 = 150 \nota{acumulado 150}
\end{aligned}$
\end{minipage}\hfill
\begin{minipage}[t]{0.49\textwidth}
\paso{Ruta 2: \ P1 $\to$ C2}
$\begin{aligned}
\text{por unidad} &= 6\\
\text{envío} &= \min(35, \infty, 20) = 20 \nota{C2 llena}\\
\text{costo} &= 6 \times 20 = 120 \nota{acumulado 270}
\end{aligned}$
\end{minipage}

\noindent
\begin{minipage}[t]{0.49\textwidth}
\paso{Ruta 3: \ P1 $\to$ C1}
$\begin{aligned}
\text{por unidad} &= 8\\
\text{envío} &= \min(15, \infty, 45) = 15 \nota{P1 vacía}\\
\text{costo} &= 8 \times 15 = 120 \nota{acumulado 390}
\end{aligned}$
\end{minipage}\hfill
\begin{minipage}[t]{0.49\textwidth}
\paso{Ruta 4: \ P2 $\to$ C1}
$\begin{aligned}
\text{por unidad} &= 9 \nota{C4 ya está llena}\\
\text{envío} &= \min(50, \infty, 30) = 30 \nota{C1 llena}\\
\text{costo} &= 9 \times 30 = 270 \nota{acumulado 660}
\end{aligned}$
\end{minipage}

\medskip
Ahora C1, C2 y C4 están completas. Solo falta C3, que pide 30; a P2 le quedan 20 y a P3 le quedan 10.

\paso{Ruta 5: \ P2 $\to$ C1 $\to$ P1 $\to$ C3 \quad (primera vez que se deshace un envío)}
Las rutas directas a C3 son caras: P2 \ra{} C3 cuesta 13 y P3 \ra{} C3 cuesta 16. Pero P1 \ra{} C3 cuesta solo 10, y
aunque a P1 ya no le queda energía, sí está mandando 15 a C1. En la red de lo que queda existe el tramo al revés
C1 \ra{} P1, con espacio 15 y costo $-8$. La ruta más barata lo aprovecha:
\begin{align*}
\text{costo por unidad} &= 9 - 8 + 10 = 11 \nota{el tramo al revés resta su costo}\\
\text{envío} &= \min(20,\ \infty,\ 15,\ \infty,\ 30) = 15 \nota{solo se pueden deshacer las 15 de P1 a C1}\\
\text{costo} &= 11 \times 15 = 165, \qquad \text{acumulado} = 825 \nota{van 110 de 125}
\end{align*}

\begin{ojo}
\textbf{¿Qué significa recorrer C1 \ra{} P1 al revés?} Que P2 le manda a C1 \textbf{15 unidades más} (cuesta 9 cada una);
como C1 ya estaba completa, P1 \textbf{deja de mandarle} esas 15 (ahorra 8 cada una); y P1 usa esa energía para
abastecer a C3 (cuesta 10 cada una). C1 sigue recibiendo 45 y C3 recibe 15. El cambio neto cuesta
$9 - 8 + 10 = 11$ por unidad, más barato que cualquier envío directo a C3.

\begin{center}
\begin{tikzpicture}[x=1.15cm, y=-1.0cm]
  \begin{scope}
    \node[fuente] (P1) at (0,0.5) {P1}; \node[fuente] (P2) at (0,2) {P2};
    \node[sumidero] (C1) at (3.4,0) {C1}; \node[sumidero] (C3) at (3.4,2.67) {C3};
    \draw[usado] (P1) -- node[etq, above] {15} (C1);
    \draw[usado] (P2) -- node[etq, pos=0.15] {30} (C1);
    \node[font=\small\bfseries] at (1.7,3.5) {Antes (lo enviado)};
  \end{scope}
  \begin{scope}[xshift=5.4cm]
    \node[fuente] (P1) at (0,0.5) {P1}; \node[fuente] (P2) at (0,2) {P2};
    \node[sumidero] (C1) at (3.4,0) {C1}; \node[sumidero] (C3) at (3.4,2.67) {C3};
    \draw[camino] (P2) -- node[etq, pos=0.15] {$+9$} (C1);
    \draw[regreso] (C1) -- node[etq, above] {$-8$} (P1);
    \draw[camino] (P1) -- node[etq, pos=0.75] {$+10$} (C3);
    \node[font=\small\bfseries] at (1.7,3.5) {Ruta 5 (15 unidades)};
  \end{scope}
  \begin{scope}[xshift=10.8cm]
    \node[fuente] (P1) at (0,0.5) {P1}; \node[fuente] (P2) at (0,2) {P2};
    \node[sumidero] (C1) at (3.4,0) {C1}; \node[sumidero] (C3) at (3.4,2.67) {C3};
    \draw[gris, dotted] (P1) -- node[etq, above, text=gray] {0} (C1);
    \draw[usado] (P2) -- node[etq, pos=0.15] {45} (C1);
    \draw[usado] (P1) -- node[etq, pos=0.75] {15} (C3);
    \node[font=\small\bfseries] at (1.7,3.5) {Después};
  \end{scope}
\end{tikzpicture}
\end{center}
El tramo punteado grueso es el que se recorre al revés. Después de la ruta 5, P1 \ra{} C1 queda en 0.
\end{ojo}

\paso{Ruta 6: \ P2 $\to$ C3}
A C3 le faltan 15; a P2 le quedan 5 y a P3, 10.
\begin{align*}
\text{costo por unidad} &= 13 \nota{empate: P3 \ra{} C2 \ra{} P1 \ra{} C3 también cuesta 13}\\
\text{envío} &= \min(5,\ \infty,\ 15) = 5 \nota{a P2 solo le quedan 5}\\
\text{costo} &= 13 \times 5 = 65, \qquad \text{acumulado} = 890 \nota{van 115 de 125}
\end{align*}
Hay dos rutas de 13 y se toma una de ellas; la otra se usa en la vuelta siguiente. P2 queda sin energía.

\paso{Ruta 7: \ P3 $\to$ C2 $\to$ P1 $\to$ C3 \quad (otra vez se deshace un envío)}
Solo a P3 le queda energía (10) y solo a C3 le falta (10). Ir directo de P3 a C3 cuesta 16; es más barato que P3
abastezca a C2 y que P1 cambie parte de lo que mandaba a C2 hacia C3.
\begin{align*}
\text{costo por unidad} &= 9 - 6 + 10 = 13 \nota{C2 \ra{} P1 deshace P1 \ra{} C2: ahorra 6}\\
\text{envío} &= \min(10,\ \infty,\ 20,\ \infty,\ 10) = 10 \nota{a P3 le quedan 10; a C3 le faltan 10}\\
\text{costo} &= 13 \times 10 = 130, \qquad \text{acumulado} = 1\,020 \nota{van 125 de 125: terminado}
\end{align*}

\begin{center}
\begin{tikzpicture}[x=1.15cm, y=-1.0cm]
  \begin{scope}
    \node[fuente] (P1) at (0,0.5) {P1}; \node[fuente] (P3) at (0,3.5) {P3};
    \node[sumidero] (C2) at (3.4,1.33) {C2}; \node[sumidero] (C3) at (3.4,2.67) {C3};
    \draw[usado] (P1) -- node[etq, pos=0.4] {20} (C2);
    \draw[usado] (P1) -- node[etq, pos=0.75] {15} (C3);
    \node[font=\small\bfseries] at (1.7,4.3) {Antes (lo enviado)};
  \end{scope}
  \begin{scope}[xshift=5.4cm]
    \node[fuente] (P1) at (0,0.5) {P1}; \node[fuente] (P3) at (0,3.5) {P3};
    \node[sumidero] (C2) at (3.4,1.33) {C2}; \node[sumidero] (C3) at (3.4,2.67) {C3};
    \draw[camino] (P3) -- node[etq, pos=0.2] {$+9$} (C2);
    \draw[regreso] (C2) -- node[etq, pos=0.4, above] {$-6$} (P1);
    \draw[camino] (P1) -- node[etq, pos=0.86] {$+10$} (C3);
    \node[font=\small\bfseries] at (1.7,4.3) {Ruta 7 (10 unidades)};
  \end{scope}
  \begin{scope}[xshift=10.8cm]
    \node[fuente] (P1) at (0,0.5) {P1}; \node[fuente] (P3) at (0,3.5) {P3};
    \node[sumidero] (C2) at (3.4,1.33) {C2}; \node[sumidero] (C3) at (3.4,2.67) {C3};
    \draw[usado] (P1) -- node[etq, pos=0.4] {10} (C2);
    \draw[usado] (P1) -- node[etq, pos=0.86] {25} (C3);
    \draw[usado] (P3) -- node[etq, pos=0.2] {10} (C2);
    \node[font=\small\bfseries] at (1.7,4.3) {Después};
  \end{scope}
\end{tikzpicture}
\end{center}

\paso{Resumen de las siete rutas}
\begin{center}\small
\begin{tabular}{clcccc}
\toprule
\textbf{Ruta} & \textbf{Recorrido} & \textbf{Costo por unidad} & \textbf{Envío} & \textbf{Costo} & \textbf{Acumulado} \\
\midrule
1 & P3 \ra{} C4 & 5 & 30 & 150 & 150 \\
2 & P1 \ra{} C2 & 6 & 20 & 120 & 270 \\
3 & P1 \ra{} C1 & 8 & 15 & 120 & 390 \\
4 & P2 \ra{} C1 & 9 & 30 & 270 & 660 \\
5 & P2 \ra{} C1 \ra{} P1 \ra{} C3 & $9 - 8 + 10 = 11$ & 15 & 165 & 825 \\
6 & P2 \ra{} C3 & 13 & 5 & 65 & 890 \\
7 & P3 \ra{} C2 \ra{} P1 \ra{} C3 & $9 - 6 + 10 = 13$ & 10 & 130 & 1\,020 \\
\bottomrule
\end{tabular}
\end{center}
Fíjate en que el costo por unidad de las rutas \textbf{nunca baja}: 5, 6, 8, 9, 11, 13, 13. Cada vuelta usa la ruta más
barata que queda, y las baratas se van agotando.

\subsection*{Resultado}
\[
\res{\text{Costo mínimo } Z^* = \$1\,020}
\]
\noindent
\begin{minipage}[c]{0.42\textwidth}\small\centering
\begin{tabular}{lccc}
\toprule
\textbf{Arco} & \textbf{Envío} & \textbf{Precio} & \textbf{Costo} \\
\midrule
P1 \ra{} C2 & 10 & 6  & 60 \\
P1 \ra{} C3 & 25 & 10 & 250 \\
P2 \ra{} C1 & 45 & 9  & 405 \\
P2 \ra{} C3 & 5  & 13 & 65 \\
P3 \ra{} C2 & 10 & 9  & 90 \\
P3 \ra{} C4 & 30 & 5  & 150 \\
\midrule
\textbf{Total} & 125 & & \textbf{1\,020} \\
\bottomrule
\end{tabular}
\end{minipage}\hfill
\begin{minipage}[c]{0.56\textwidth}\small
Los otros seis arcos llevan 0. Es la misma solución óptima que reporta Winston ($x_{12} = 10$, $x_{13} = 25$,
$x_{21} = 45$, $x_{23} = 5$, $x_{32} = 10$, $x_{34} = 30$).

\textbf{Comprobación (conservación en cada nodo):}
\begin{itemize}[itemsep=1pt, leftmargin=*]
  \item Sale: P1 $10 + 25 = 35$, P2 $45 + 5 = 50$, P3 $10 + 30 = 40$ \checkmark
  \item Entra: C1 $45$, C2 $10 + 10 = 20$, C3 $25 + 5 = 30$, C4 $30$ \checkmark
  \item Capacidades: ningún arco tiene límite \checkmark
  \item Costo: $60 + 250 + 405 + 65 + 90 + 150 = 1\,020$ \checkmark
\end{itemize}
\end{minipage}

\subsection*{Interpretación gráfica}
\noindent
\begin{minipage}[c]{0.46\textwidth}\centering
\begin{tikzpicture}[x=1.5cm, y=-1.25cm]
  \node[fuente, label=left:{\small $+35$}] (P1) at (0,0.5) {P1};
  \node[fuente, label=left:{\small $+50$}] (P2) at (0,2) {P2};
  \node[fuente, label=left:{\small $+40$}] (P3) at (0,3.5) {P3};
  \node[sumidero, label=right:{\small $-45$}] (C1) at (3.4,0) {C1};
  \node[sumidero, label=right:{\small $-20$}] (C2) at (3.4,1.33) {C2};
  \node[sumidero, label=right:{\small $-30$}] (C3) at (3.4,2.67) {C3};
  \node[sumidero, label=right:{\small $-30$}] (C4) at (3.4,4) {C4};
  \foreach \p/\c in {P1/C1, P1/C4, P2/C2, P2/C4, P3/C1, P3/C3} \draw[gris] (\p) -- (\c);
  \foreach \p/\c/\f in {P1/C2/10, P1/C3/25, P2/C1/45, P2/C3/5, P3/C2/10, P3/C4/30}
    \draw[usado] (\p) -- node[etq, pos=0.8, font=\scriptsize\bfseries] {\f} (\c);
\end{tikzpicture}
\end{minipage}\hfill
\begin{minipage}[c]{0.51\textwidth}
Las flechas gruesas son los envíos óptimos (con las unidades); las grises, los arcos que no se usan.

\textbf{¿Qué significa?} La planta 3, que está cerca de la ciudad 4 (costo 5), la abastece completa y le da el resto a
la ciudad 2. La planta 2 se dedica casi por completo a la ciudad 1. La planta 1 reparte entre las ciudades 2 y 3. La
ciudad 3 es la más cara de abastecer desde cualquier planta, por eso queda para el final y el método tuvo que
\textbf{reacomodar} dos veces los envíos de la planta 1 para atenderla al menor costo. Con 7 rutas se encontró el plan
de \textbf{\$1\,020}.
\end{minipage}

% =====================================================================
\section{Distribution Unlimited Co.: transbordo con capacidades}
\textit{Fuente: Hillier y Lieberman, Introducción a la Investigación de Operaciones, capítulo de modelos de optimización
de redes (problema del flujo de costo mínimo).}

\subsection*{Planteamiento del problema}
La Distribution Unlimited Co. fabrica un producto en \textbf{dos fábricas} y debe llevarlo a \textbf{dos almacenes}.
La fábrica F1 produce 50 unidades y la F2 produce 40; el almacén A1 necesita 30 y el A2 necesita 60. La mercadería
puede ir directo, pasar por un \textbf{centro de distribución} (CD) o moverse entre los almacenes. Dos tramos tienen
límite: el camión de F1 a F2 lleva como mucho 10 unidades y el de CD a A2, como mucho 80.

La pregunta: \textbf{¿cuánto mandar por cada arco para abastecer los dos almacenes al menor costo?}

\noindent
\begin{minipage}[c]{0.40\textwidth}\small\centering
\begin{tabular}{lcc}
\toprule
\textbf{Arco} & \textbf{Costo} & \textbf{Capacidad} \\
\midrule
F1 \ra{} F2 & \$200 & 10 \\
F1 \ra{} CD & \$400 & sin límite \\
F1 \ra{} A1 & \$900 & sin límite \\
F2 \ra{} CD & \$300 & sin límite \\
CD \ra{} A2 & \$100 & 80 \\
A1 \ra{} A2 & \$300 & sin límite \\
A2 \ra{} A1 & \$200 & sin límite \\
\bottomrule
\end{tabular}
\end{minipage}\hfill
\begin{minipage}[c]{0.57\textwidth}\centering
\begin{tikzpicture}[x=1.45cm, y=-1.0cm]
  \node[fuente, label=left:{\small $+50$}] (F1) at (0,0.4) {F1};
  \node[fuente, label=left:{\small $+40$}] (F2) at (0,3.2) {F2};
  \node[nodo, label=below:{\small $0$}] (CD) at (2,1.8) {CD};
  \node[sumidero, label=right:{\small $-30$}] (A1) at (4,0.4) {A1};
  \node[sumidero, label=right:{\small $-60$}] (A2) at (4,3.2) {A2};
  \draw[arco] (F1) -- node[etq, align=center] {200\\(máx.\ 10)} (F2);
  \draw[arco] (F1) -- node[etq] {400} (CD);
  \draw[arco] (F1) -- node[etq, above] {900} (A1);
  \draw[arco] (F2) -- node[etq] {300} (CD);
  \draw[arco] (CD) -- node[etq, align=center] {100\\(máx.\ 80)} (A2);
  \draw[arco] (A1) to[bend left=22] node[etq, right] {300} (A2);
  \draw[arco] (A2) to[bend left=22] node[etq, left] {200} (A1);
\end{tikzpicture}
\end{minipage}

\subsection*{Formulación del modelo}
Sean $\x{F1}{CD}$, $\x{F1}{A1}$, \dots{} las unidades que van por cada arco (7 variables). \textbf{Conservación}
(sale $-$ entra $=$ oferta): F1 ofrece 50, F2 ofrece 40, CD es de transbordo (0), A1 pide 30 ($-30$) y A2 pide 60
($-60$). \textbf{Capacidades:} $\x{F1}{F2} \le 10$ y $\x{CD}{A2} \le 80$.
\[
\text{Min } Z = 200\,\x{F1}{F2} + 400\,\x{F1}{CD} + 900\,\x{F1}{A1} + 300\,\x{F2}{CD} + 100\,\x{CD}{A2}
                  + 300\,\x{A1}{A2} + 200\,\x{A2}{A1}
\]
\[
\begin{aligned}
\text{s.a.}\quad & \x{F1}{F2} + \x{F1}{CD} + \x{F1}{A1} = 50 \nota{F1} &\qquad& \x{A1}{A2} - \x{F1}{A1} - \x{A2}{A1} = -30 \nota{A1}\\
& \x{F2}{CD} - \x{F1}{F2} = 40 \nota{F2} && \x{A2}{A1} - \x{CD}{A2} - \x{A1}{A2} = -60 \nota{A2}\\
& \x{CD}{A2} - \x{F1}{CD} - \x{F2}{CD} = 0 \nota{CD} && \x{F1}{F2} \le 10, \ \ \x{CD}{A2} \le 80, \ \ \text{todas} \ge 0
\end{aligned}
\]

\subsection*{Preparación}
Se ofrecen $50 + 40 = 90$ y se piden $30 + 60 = 90$: \textbf{está balanceado}. \Sini{} se une a F1 (espacio 50) y a
F2 (espacio 40); A1 (espacio 30) y A2 (espacio 60) se unen a \Tfin. CD no ofrece ni pide: es un nodo de paso. Al inicio
todo está vacío, costo acumulado 0.

% ---------------------------------------------------------------------
\subsection*{Las rutas, una por una}
En los dibujos, la ruta de cada vuelta va en trazo grueso con el costo de cada tramo.

\newcommand{\nodosDU}{%
  \node[fuente, font=\scriptsize\bfseries] (F1) at (0,0.4) {F1};
  \node[fuente, font=\scriptsize\bfseries] (F2) at (0,3.2) {F2};
  \node[nodo, font=\scriptsize\bfseries] (CD) at (2,1.8) {CD};
  \node[sumidero, font=\scriptsize\bfseries] (A1) at (4,0.4) {A1};
  \node[sumidero, font=\scriptsize\bfseries] (A2) at (4,3.2) {A2};}
\newcommand{\grisDU}{%
  \draw[gris] (F1) -- (F2); \draw[gris] (F1) -- (CD); \draw[gris] (F1) -- (A1); \draw[gris] (F2) -- (CD);
  \draw[gris] (CD) -- (A2); \draw[gris] (A1) to[bend left=22] (A2); \draw[gris] (A2) to[bend left=22] (A1);}

\paso{Ruta 1: \ F2 $\to$ CD $\to$ A2}
\noindent
\begin{minipage}[c]{0.32\textwidth}\centering
\begin{tikzpicture}[x=1.1cm, y=-0.8cm]
  \nodosDU \grisDU
  \draw[camino] (F2) -- node[etq] {300} (CD);
  \draw[camino] (CD) -- node[etq] {100} (A2);
\end{tikzpicture}
\end{minipage}\hfill
\begin{minipage}[c]{0.66\textwidth}
\begin{align*}
\text{por unidad} &= 300 + 100 = 400 \nota{la más barata}\\
\text{envío} &= \min(40, \infty, 80, 60) = 40 \nota{a F2 le quedan 40}\\
\text{costo} &= 400 \times 40 = 16\,000\\
\text{acumulado} &= 16\,000 \nota{van 40 de 90}
\end{align*}
\end{minipage}

F2 manda toda su producción al almacén 2 pasando por el centro. En CD \ra{} A2 quedan $80 - 40 = 40$ de espacio.

\paso{Ruta 2: \ F1 $\to$ CD $\to$ A2}
\noindent
\begin{minipage}[c]{0.32\textwidth}\centering
\begin{tikzpicture}[x=1.1cm, y=-0.8cm]
  \nodosDU \grisDU
  \draw[camino] (F1) -- node[etq] {400} (CD);
  \draw[camino] (CD) -- node[etq] {100} (A2);
\end{tikzpicture}
\end{minipage}\hfill
\begin{minipage}[c]{0.66\textwidth}
\begin{align*}
\text{por unidad} &= 400 + 100 = 500 \nota{pasando por F2: 600}\\
\text{envío} &= \min(50, \infty, 40, 20) = 20 \nota{a A2 le faltan 20}\\
\text{costo} &= 500 \times 20 = 10\,000\\
\text{acumulado} &= 26\,000 \nota{van 60 de 90}
\end{align*}
\end{minipage}

A2 queda completo (60). En CD \ra{} A2 van 60 de 80: quedan 20 de espacio. A F1 le quedan 30.

\paso{Ruta 3: \ F1 $\to$ CD $\to$ A2 $\to$ A1}
\noindent
\begin{minipage}[c]{0.32\textwidth}\centering
\begin{tikzpicture}[x=1.1cm, y=-0.8cm]
  \nodosDU \grisDU
  \draw[camino] (F1) -- node[etq] {400} (CD);
  \draw[camino] (CD) -- node[etq] {100} (A2);
  \draw[camino] (A2) to[bend left=22] node[etq, left] {200} (A1);
\end{tikzpicture}
\end{minipage}\hfill
\begin{minipage}[c]{0.66\textwidth}
\begin{align*}
\text{por unidad} &= 400 + 100 + 200 = 700 \nota{directo: 900}\\
\text{envío} &= \min(30, \infty, 20, \infty, 30) = 20 \nota{límite de CD \ra{} A2}\\
\text{costo} &= 700 \times 20 = 14\,000\\
\text{acumulado} &= 40\,000 \nota{van 80 de 90}
\end{align*}
\end{minipage}

Ahora solo falta A1. Es más barato llegar a él \textbf{pasando por A2} que ir directo, hasta que el tramo CD \ra{} A2
se llena: lleva $40 + 20 + 20 = 80$, su capacidad.

\paso{Ruta 4: \ F1 $\to$ A1}
\noindent
\begin{minipage}[c]{0.32\textwidth}\centering
\begin{tikzpicture}[x=1.1cm, y=-0.8cm]
  \nodosDU \grisDU
  \draw[corte, gray] (CD) -- node[etq, text=black, font=\tiny] {lleno} (A2);
  \draw[camino] (F1) -- node[etq, above] {900} (A1);
\end{tikzpicture}
\end{minipage}\hfill
\begin{minipage}[c]{0.66\textwidth}
\begin{align*}
\text{por unidad} &= 900 \nota{CD \ra{} A2 ya está lleno}\\
\text{envío} &= \min(10, \infty, 10) = 10 \nota{a A1 le faltan 10}\\
\text{costo} &= 900 \times 10 = 9\,000\\
\text{acumulado} &= 49\,000 \nota{90 de 90: fin}
\end{align*}
\end{minipage}

En este ejercicio ninguna ruta tuvo que deshacer envíos: lo que marcó el camino fue la \textbf{capacidad} de 80 del
tramo CD \ra{} A2, que obligó a usar al final el arco caro F1 \ra{} A1.

\subsection*{Resultado}
\[
\res{\text{Costo mínimo } Z^* = \$49\,000}
\]
\noindent
\begin{minipage}[c]{0.42\textwidth}\small\centering
\begin{tabular}{lccc}
\toprule
\textbf{Arco} & \textbf{Envío} & \textbf{Precio} & \textbf{Costo} \\
\midrule
F1 \ra{} CD & 40 & 400 & 16\,000 \\
F1 \ra{} A1 & 10 & 900 & 9\,000 \\
F2 \ra{} CD & 40 & 300 & 12\,000 \\
CD \ra{} A2 & 80 & 100 & 8\,000 \\
A2 \ra{} A1 & 20 & 200 & 4\,000 \\
\midrule
\textbf{Total} & & & \textbf{49\,000} \\
\bottomrule
\end{tabular}
\end{minipage}\hfill
\begin{minipage}[c]{0.56\textwidth}\small
F1 \ra{} F2 y A1 \ra{} A2 llevan 0. Coincide con la solución de Hillier y Lieberman.

\textbf{Comprobación} (sale $-$ entra $=$ oferta, y capacidades):
\begin{itemize}[itemsep=1pt, leftmargin=*]
  \item F1: $50 - 0 = 50$ \checkmark \quad F2: $40 - 0 = 40$ \checkmark
  \item CD: $80 - (40 + 40) = 0$ \checkmark \ (nodo de paso)
  \item A1: $0 - (10 + 20) = -30$ \checkmark \quad A2: $20 - 80 = -60$ \checkmark
  \item F1 \ra{} F2: $0 \le 10$ \checkmark \quad CD \ra{} A2: $80 \le 80$ \checkmark (al tope)
  \item Costo: $16\,000 + 9\,000 + 12\,000 + 8\,000 + 4\,000 = 49\,000$ \checkmark
\end{itemize}
\end{minipage}

\subsection*{Interpretación gráfica}
\noindent
\begin{minipage}[c]{0.47\textwidth}\centering
\begin{tikzpicture}[x=1.3cm, y=-1.0cm]
  \node[fuente, label=left:{\small $+50$}] (F1) at (0,0.4) {F1};
  \node[fuente, label=left:{\small $+40$}] (F2) at (0,3.2) {F2};
  \node[nodo, label=below:{\small $0$}] (CD) at (2,1.8) {CD};
  \node[sumidero, label=right:{\small $-30$}] (A1) at (4,0.4) {A1};
  \node[sumidero, label=right:{\small $-60$}] (A2) at (4,3.2) {A2};
  \draw[gris] (F1) -- (F2);
  \draw[gris] (A1) to[bend left=22] (A2);
  \draw[usado] (F1) -- node[etq, font=\scriptsize\bfseries] {40} (CD);
  \draw[usado] (F1) -- node[etq, above, font=\scriptsize\bfseries] {10} (A1);
  \draw[usado] (F2) -- node[etq, font=\scriptsize\bfseries] {40} (CD);
  \draw[usado] (CD) -- node[etq, font=\scriptsize\bfseries] {80 (lleno)} (A2);
  \draw[usado] (A2) to[bend left=22] node[etq, left, font=\scriptsize\bfseries] {20} (A1);
\end{tikzpicture}
\end{minipage}\hfill
\begin{minipage}[c]{0.50\textwidth}
\textbf{¿Qué significa?} Casi todo pasa por el centro de distribución: las 40 unidades de F2 y 40 de las 50 de F1
llegan a CD y de ahí 80 van al almacén 2, que es todo lo que permite ese tramo. El almacén 2 se queda con 60 y
\textbf{reenvía 20 al almacén 1}, porque eso ($400 + 100 + 200 = 700$) sale más barato que el envío directo (900). Las
últimas 10 unidades de F1 van directo al almacén 1. El camión entre fábricas no se usa. Costo: \textbf{\$49\,000}.
\end{minipage}

% =====================================================================
\section{Ruta más corta en el parque Seervada}
\textit{Fuente: Hillier y Lieberman, Introducción a la Investigación de Operaciones, capítulo de modelos de optimización
de redes (problema de la ruta más corta).}

\subsection*{Planteamiento del problema}
El parque Seervada tiene caminos de tierra entre la entrada (O), varias estaciones (A, B, C, D y E) y un mirador (T).
Cada camino se puede recorrer en los dos sentidos. Las distancias, en millas, son:

\noindent
\begin{minipage}[c]{0.38\textwidth}\small\centering
\begin{tabular}{lc@{\quad}lc}
\toprule
\textbf{Camino} & \textbf{Millas} & \textbf{Camino} & \textbf{Millas} \\
\midrule
O -- A & 2 & B -- D & 4 \\
O -- B & 5 & B -- E & 3 \\
O -- C & 4 & C -- E & 4 \\
A -- B & 2 & D -- E & 1 \\
A -- D & 7 & D -- T & 5 \\
B -- C & 1 & E -- T & 7 \\
\bottomrule
\end{tabular}
\end{minipage}\hfill
\begin{minipage}[c]{0.60\textwidth}\centering
\begin{tikzpicture}[x=1.45cm, y=-0.95cm]
  \node[fuente, label=left:{\small $+1$}] (O) at (0,2) {O};
  \node[nodo] (A) at (1.5,0.6) {A};
  \node[nodo] (B) at (2,2) {B};
  \node[nodo] (C) at (1.5,3.4) {C};
  \node[nodo] (D) at (3.6,0.9) {D};
  \node[nodo] (E) at (3.6,3.1) {E};
  \node[sumidero, label=right:{\small $-1$}] (T) at (5,2) {T};
  \foreach \u/\v/\d in {O/A/2, O/B/5, O/C/4, A/B/2, A/D/7, B/C/1, B/D/4, B/E/3, C/E/4, D/E/1, D/T/5, E/T/7}
    \draw[calle] (\u) -- node[etq] {\d} (\v);
\end{tikzpicture}
\end{minipage}

La pregunta: \textbf{¿cuál es la ruta más corta de la entrada O al mirador T?}

\subsection*{Formulación del modelo}
La ruta más corta es un flujo a costo mínimo con datos especiales: la entrada O \textbf{ofrece 1 unidad} ($+1$), el
mirador T \textbf{la pide} ($-1$) y las demás estaciones son de paso (0); esa unidad es ``el viajero''. Cada camino se
convierte en \textbf{dos arcos} (ida y vuelta) con el mismo costo, su distancia: son 24 arcos, ninguno con límite. Así,
el costo total es la longitud de la ruta por donde pasa la unidad.
\[
\begin{aligned}
\text{Min } Z &= 2\,\x{O}{A} + 5\,\x{O}{B} + 4\,\x{O}{C} + 2\,\x{A}{B} + \dots + 5\,\x{T}{D} + 7\,\x{T}{E} \nota{24 arcos}\\
\text{s.a.}\quad & \text{O: sale} - \text{entra} = 1, \qquad \text{T: sale} - \text{entra} = -1\\
& \text{A, B, C, D, E: sale} - \text{entra} = 0 \nota{son puntos de paso}\\
& \text{todas} \ge 0
\end{aligned}
\]
Por ejemplo, en A la regla es $\x{A}{B} + \x{A}{D} + \x{A}{O} - \x{O}{A} - \x{B}{A} - \x{D}{A} = 0$.
Gracias a la propiedad de soluciones enteras, cada arco termina llevando 0 o 1: la unidad no se ``parte'' entre dos caminos.

\subsection*{Preparación}
Se ofrece 1 y se pide 1: \textbf{está balanceado}. \Sini{} se une a O con espacio 1 y T se une a \Tfin{} con espacio 1.
Todas las distancias son positivas, así que no puede haber recorridos circulares de costo negativo.

% ---------------------------------------------------------------------
\subsection*{La única ruta}

\paso{Ruta 1: \ O $\to$ A $\to$ B $\to$ D $\to$ T}
Para encontrarla se anota, para cada estación, \textbf{la distancia más corta para llegar desde O} (es lo que hace la
búsqueda de la ruta más barata). Se empieza con O en 0 y se va corrigiendo hasta que nada cambie:

\begin{center}\small
\begin{tabular}{cccp{6.4cm}}
\toprule
\textbf{Estación} & \textbf{Distancia desde O} & \textbf{Se llega desde} & \textbf{Por qué} \\
\midrule
O & 0  & --- & es el inicio \\
A & 2  & O & O -- A: 2 \\
B & 4  & A & $2 + 2 = 4$, mejor que O -- B directo (5) \\
C & 4  & O & O -- C: 4, mejor que por B ($4 + 1 = 5$) \\
E & 7  & B & $4 + 3 = 7$, mejor que por C ($4 + 4 = 8$) \\
D & 8  & B & $4 + 4 = 8$; por E también da $7 + 1 = 8$ (empate); por A, $2 + 7 = 9$ \\
T & 13 & D & $8 + 5 = 13$, mejor que por E ($7 + 7 = 14$) \\
\bottomrule
\end{tabular}
\end{center}

Siguiendo la columna ``se llega desde'' hacia atrás, desde T: T $\leftarrow$ D $\leftarrow$ B $\leftarrow$ A $\leftarrow$ O.
\begin{align*}
\text{costo por unidad} &= 2 + 2 + 4 + 5 = 13 \nota{la suma de las distancias}\\
\text{envío} &= \min(1,\ \infty,\ \infty,\ \infty,\ \infty,\ 1) = 1 \nota{solo hay un viajero}\\
\text{costo} &= 13 \times 1 = 13, \qquad \text{acumulado} = 13 \nota{va 1 de 1: terminado}
\end{align*}

Con una sola vuelta se entregó todo lo pedido: no hace falta otra iteración. En la ruta más corta el método
siempre termina en una vuelta, porque solo hay una unidad que mandar.

\subsection*{Resultado}
\[
\res{\text{Distancia mínima de O a T} = 13 \text{ millas}}
\]
La ruta es \textbf{O \ra{} A \ra{} B \ra{} D \ra{} T}: los arcos O \ra{} A, A \ra{} B, B \ra{} D y D \ra{} T llevan 1
unidad y los otros 20 llevan 0.

\textbf{Hay más de una ruta óptima.} Hillier y Lieberman reportan dos rutas de 13 millas:
O \ra{} A \ra{} B \ra{} D \ra{} T y O \ra{} A \ra{} B \ra{} E \ra{} D \ra{} T ($2 + 2 + 3 + 1 + 5 = 13$). El empate se ve en la
tabla (a D se llega con 8 millas desde B y desde E); la calculadora encuentra una de ellas.

\textbf{Comprobación (conservación en cada nodo):}
\begin{itemize}[itemsep=2pt]
  \item O: sale 1, entra 0: $1 - 0 = 1$ \checkmark \qquad T: sale 0, entra 1: $0 - 1 = -1$ \checkmark
  \item A, B y D: entra 1 y sale 1: $1 - 1 = 0$ \checkmark \qquad C y E: no pasa nada, $0 - 0 = 0$ \checkmark
  \item Capacidades: no hay límites \checkmark. \ Distancia: $2 + 2 + 4 + 5 = 13$ \checkmark
\end{itemize}

\subsection*{Interpretación gráfica}
\noindent
\begin{minipage}[c]{0.56\textwidth}\centering
\begin{tikzpicture}[x=1.4cm, y=-1.0cm]
  \node[fuente, label=left:{\small $+1$}] (O) at (0,2) {O};
  \node[nodo, label=above:{\scriptsize 2}] (A) at (1.5,0.6) {A};
  \node[nodo, label=below:{\scriptsize 4}] (B) at (2,2) {B};
  \node[nodo, label=below:{\scriptsize 4}] (C) at (1.5,3.4) {C};
  \node[nodo, label=above:{\scriptsize 8}] (D) at (3.6,0.9) {D};
  \node[nodo, label=below:{\scriptsize 7}] (E) at (3.6,3.1) {E};
  \node[sumidero, label=right:{\small $-1$}, label=above:{\scriptsize 13}] (T) at (5,2) {T};
  \foreach \u/\v/\d in {O/B/5, O/C/4, A/D/7, B/C/1, C/E/4, E/T/7}
    \draw[gray!55] (\u) -- node[etq, text=gray] {\d} (\v);
  \draw[very thick, dash pattern=on 3pt off 2pt] (B) -- node[etq] {3} (E);
  \draw[very thick, dash pattern=on 3pt off 2pt] (E) -- node[etq] {1} (D);
  \draw[camino] (O) -- node[etq] {2} (A);
  \draw[camino] (A) -- node[etq] {2} (B);
  \draw[camino] (B) -- node[etq] {4} (D);
  \draw[camino] (D) -- node[etq] {5} (T);
\end{tikzpicture}
\end{minipage}\hfill
\begin{minipage}[c]{0.42\textwidth}
La ruta gruesa es la que encontró el método; el tramo B -- E -- D en trazo discontinuo es la otra ruta de 13 millas.
El número pequeño junto a cada estación es su distancia más corta desde O.
\end{minipage}

\textbf{¿Qué significa?} Para ir de la entrada al mirador conviene pasar por A, B y D: son \textbf{13 millas}, y
ningún otro recorrido es más corto (por ejemplo, O \ra{} B \ra{} D \ra{} T mide 14 y O \ra{} A \ra{} D \ra{} T mide 14).
El ejercicio muestra que la ruta más corta no necesita un método aparte: es un flujo a costo mínimo en el que se manda
una sola unidad.

% =====================================================================
\section*{Resumen de los tres ejercicios}
\addcontentsline{toc}{part}{Resumen y conclusiones}
\begin{center}\small
\begin{tabular}{llccll}
\toprule
\textbf{Ejercicio} & \textbf{Caso particular} & \textbf{Nodos / arcos} & \textbf{Rutas} & \textbf{¿Deshace?} & \textbf{Óptimo} \\
\midrule
Powerco & transporte & 7 / 12 & 7 & sí (rutas 5 y 7) & \$1\,020 \\
Distribution Unl. & transbordo y límites & 5 / 7 & 4 & no & \$49\,000 \\
Seervada & ruta más corta & 7 / 24 & 1 & no & 13 millas \\
\bottomrule
\end{tabular}
\end{center}

\begin{idea}[Conclusiones]
\begin{itemize}[itemsep=2pt, leftmargin=*]
  \item El flujo a costo mínimo pide entregar todo lo pedido gastando lo menos, respetando la \textbf{conservación}
        (sale $-$ entra $=$ oferta) y las \textbf{capacidades}.
  \item Es el problema ``padre'': el transporte, la asignación, el transbordo, la ruta más corta y el flujo máximo son
        casos particulares, y un mismo método los resuelve.
  \item El método de \textbf{rutas más baratas sucesivas} junta las ofertas en \Sini{} y las demandas en \Tfin, y en cada
        vuelta manda lo que pueda por la ruta más barata con espacio.
  \item Los \textbf{tramos al revés} (costo negativo) permiten deshacer envíos y corregir el plan; por eso se busca la
        ruta con Bellman-Ford. Termina cuando se entregó todo, y con datos enteros los envíos salen enteros.
\end{itemize}
\end{idea}

\textbf{Verificación:} los tres ejercicios se comprobaron con la implementación del método (\emph{Simplex Arcade},
modo 5P Costo mín., memorias 13, 14 y 15), que reproduce las mismas rutas y costos; además, cada problema se planteó como
programación lineal y se resolvió con el Simplex Revisado de la calculadora y con la librería javascript-lp-solver, que
confirman los tres óptimos: 1\,020, 49\,000 y 13.

\section*{Referencias}
\addcontentsline{toc}{part}{Referencias}
\begin{itemize}[leftmargin=*, itemsep=1pt]
  \item Ahuja, R. K., Magnanti, T. L. y Orlin, J. B. (1993). \textit{Network Flows: Theory, Algorithms, and Applications}. Prentice Hall.
  \item Bazaraa, M. S., Jarvis, J. J. y Sherali, H. D. \textit{Programación lineal y flujo en redes}. Limusa.
  \item Hillier, F. S. y Lieberman, G. J. \textit{Introducción a la Investigación de Operaciones}. McGraw-Hill.
        Capítulo de modelos de optimización de redes.
  \item Taha, H. A. \textit{Investigación de Operaciones}. Pearson.
  \item Winston, W. L. \textit{Investigación de operaciones: aplicaciones y algoritmos}. Thomson. Capítulo del modelo de transporte.
\end{itemize}

\end{document}
